\subsection{Continuous event-level measurement uncertainty}\label{sec:continuous_measurements}
We analyse both perfect and uncertain measurements for ten independent
catalogues at each of 50 and 500 expected detections. The injection is
$\lambda_*=(-0.3,0.35,0.01,4.7)$ and generation and inference use the same
fixed 256-million-system COMPAS population. For each observation time we draw
$N\sim\mathrm{Poisson}(T\sum_b r_b(\lambda_*))$, assign each event to a bin
with probability $P_b(\lambda_*)$, and draw its latent mass and redshift
uniformly in linear coordinates inside that bin. Each catalogue supplies a
paired perfect-measurement and uncertain-measurement dataset with identical
latent events and event counts.

To generate uncertain measurements, we draw an independent SNR for each event
from $p(\rho)\propto\rho^{-4}$ for $\rho\geq12$. With
$x_i=(\mathcal M_{c,i},z_i)$, the measurement model is
\begin{align}
 y_i&=\ln x_i+s_i\odot\eta_i,\qquad\eta_i\sim\mathcal N_2(0,I),\\
 s_i&=(0.03,0.3)\,12/\rho_i.
\end{align}
The logarithm and multiplication apply coordinate by coordinate. Thus each
measured centre is randomly displaced from the true value. The data are
$\mathcal D=\{(y_i,\rho_i)\}$, with event likelihood
$\mathcal L_i(y_i\mid x)=\mathcal N_2(y_i\mid\ln x,\mathrm{diag}(s_i^2))$.
This model follows the SNR scaling of \citet[Section 3.2.2]{Riley:2023:ApJ},
but uses Gaussian errors in log coordinates rather than their fractional
linear-coordinate errors. Our two analysis priors below also differ from
their component-mass-derived mass prior and bounded $z^2$ prior.
This is a controlled measurement model conditional on detection, not a
strain-data parameter-estimation calculation. No observed values are clipped
to the latent population support.

We generate mock event posterior samples from Bayes' theorem,
\begin{equation}
 p_{\rm PE}(x\mid y_i,\rho_i)\propto
 \mathcal L_i(y_i\mid x)\,\pi_{\rm PE}(x).
\end{equation}
For every event in the first small catalogue, containing 56 detections, we
use two normalized analysis priors on the rectangle spanning all populated
bins. Their linear-coordinate densities are constant and proportional to
$\mathcal M_c z^2$, respectively. More generally, for
$\pi_{\rm PE}(x)\propto\prod_j x_j^{h_j}$, exact independent posterior draws
are obtained from
\begin{equation}
 \ln x_j\sim\mathcal N\!\left(y_{ij}+(h_j+1)s_{ij}^2,s_{ij}^2\right),
\end{equation}
truncated to the logarithms of the prior bounds. We use $h=(0,0)$ and
$h=(1,2)$ and generate 8192 and 32768 samples per event. The extra $+1$
arises from converting the linear-coordinate prior density to log coordinates.

Posterior samples recover the event coefficients through inverse-prior weighting,
\begin{equation}
 C_{ib}\propto\frac{1}{A_b K}\sum_{k=1}^{K}
 \frac{\mathbf 1(x_{ik}\in B_b)}{\pi_{\rm PE}(x_{ik})}.
\end{equation}
The omitted event evidence is independent of $\lambda$. Samples outside
populated bins remain in the Monte Carlo denominator and contribute zero
population density. Both priors recover the analytic population posterior
within the numerical tolerances at 32768 samples per event.

For the paired GP acquisition comparison, we evaluate $C_{ib}$ by analytic
within-bin integration of the specified event likelihood. This resolves
measurement widths narrower than a population bin and separates measurement
uncertainty from finite-posterior-sample integration error. The mock samples
are a separate check of prior removal, not the training likelihood for every
catalogue. All acquisition settings and accuracy criteria are unchanged.
This fixed-injection pilot is not a coverage ensemble over different true
hyperparameters.

